\section{Absolute Fact Inclusion and Selection Interpretation}
\label{app:absolute_coverage}

\textit{Selection} measures favorable-minus-adverse inclusion within a response;
the reported paired effect is the change in this imbalance. A small effect can
therefore coexist with balanced omissions. We separately computed absolute
inclusion from the existing sentence--fact matches for all 3,840 outputs.
Table~\ref{tab:absolute_coverage} reports every model and both conditions.
Favorable fact inclusion falls from 99.31\% to 98.85\% and adverse inclusion
from 99.60\% to 98.52\%: declines of 0.45 and 1.08 percentage points,
respectively. All six facts are included in 1,858/1,920 neutral responses
(96.77\%) and 1,783/1,920 goal-conditioned responses (92.86\%).

Goal-conditioned inclusion ranges from 98.12--99.79\% for favorable facts and
95.83--99.58\% for adverse facts. Llama-3.1-8B has the largest adverse-fact
decline, from 99.17\% to 95.83\%. These measurements support high overall
source-fact retention, not complete retention or an absence of selective
omission. They depend on the fact matcher and do not establish that every
additional generated claim is supported.

\begin{table*}[t]
\centering
\small
\setlength{\tabcolsep}{3pt}
\begin{tabular}{lrrrrrr}
\toprule
Model & $F^+$ N (\%) & $F^+$ G (\%) & $F^-$ N (\%) & $F^-$ G (\%) & All six N & All six G \\
\midrule
GPT-5.4 & 99.79 & 99.79 & 99.58 & 99.58 & 157 & 157 \\
GPT-5.4 Mini & 100.00 & 99.17 & 99.79 & 99.38 & 159 & 153 \\
DeepSeek v4 & 98.96 & 98.75 & 99.58 & 99.38 & 153 & 153 \\
DeepSeek v4 (Thinking) & 99.17 & 98.12 & 99.38 & 98.12 & 153 & 145 \\
Llama 3.1 70B & 98.54 & 98.75 & 99.58 & 97.71 & 152 & 145 \\
Llama 3.1 8B & 98.96 & 98.75 & 99.17 & 95.83 & 151 & 137 \\
Qwen3 14B & 99.58 & 99.38 & 99.58 & 99.17 & 156 & 153 \\
Qwen3 14B (Thinking) & 99.38 & 98.33 & 99.79 & 99.17 & 156 & 148 \\
Qwen3 32B & 98.96 & 98.54 & 100.00 & 98.75 & 155 & 149 \\
Qwen3 32B (Thinking) & 99.79 & 98.96 & 99.58 & 99.17 & 157 & 151 \\
Qwen3 8B & 99.17 & 98.96 & 99.58 & 97.92 & 154 & 146 \\
Qwen3 8B (Thinking) & 99.38 & 98.75 & 99.58 & 98.12 & 155 & 146 \\
All models & 99.31 & 98.85 & 99.60 & 98.52 & 1,858 & 1,783 \\
\bottomrule
\end{tabular}
\caption{Absolute inclusion of supplied facts from sentence--fact matches. Each model has 160 responses per condition and 480 supplied facts of each polarity per condition. Aggregate fact percentages use 5,760 facts per polarity and condition; all-six counts use 1,920 responses per condition.}
\label{tab:absolute_coverage}
\end{table*}


\section{Ordering Robustness and Eligible Comparisons}
\label{app:ordering_robustness}

The revised \textit{ordering} definition compares exactly the same
favorable--adverse fact pairs in the neutral and goal-conditioned responses
(Equation~\ref{eq:shared_pair_ordering}). Each fact's first matched sentence is
used, including mixed-polarity sentences. A same-sentence first mention is a
tie, scored 0.5 in each arm; a missing pair is excluded from both arms.
This avoids attributing favorable ordering to a pair whose relative order
cannot be observed. It also avoids comparing different pair identities
within a scenario. Scores are averaged within scenarios and then across
eligible scenarios with equal weight.

There is no shared comparable pair for DeepSeek v4 Thinking on MED\_18 or
Llama-3.1-8B on MED\_14. These two models use 159 scenarios; the other ten use
160. Exclusions apply only to the corresponding model's ordering analysis,
not globally to these scenarios or to other dimensions. Shared-pair effects
are positive for all 12 models (0.0354--0.1705), with 11 significant under the
main 60-test correction. GPT-5.4 Mini is positive but uncertain (0.0354;
95\% interval [−0.0028, 0.0747]; adjusted $q=0.1295$).

As a separate robustness check, we retain only scenarios where all six facts
are matched in both responses. Every one of the nine favorable--adverse pairs
is then shared. The same scoring rule gives positive means for all models
(0.0241--0.1734) and significance for 11 after correction across these 12
subset tests. GPT-5.4 Mini is again the exception (0.0241; 95\% interval
[−0.0157, 0.0636]). Table~\ref{tab:ordering_robustness} reports all results and
the 132--154 eligible scenarios per model.

Both analyses condition on generated content. They measure order among
jointly retained facts, not the order of missing evidence. Selective omission
can still change eligibility and therefore the conditional estimate. The
all-six-facts subset is a sensitivity analysis, not an estimate for all
responses. This revision changes both mixed-sentence handling and tie
treatment, so it need not equal the original score even on complete outputs.

\begin{table*}[t]
\centering
\scriptsize
\setlength{\tabcolsep}{3pt}
\begin{tabular}{lrrlrrl}
\toprule
Model & $n$ & Shared-pair & 95\% CI & $n_6$ & All-six & 95\% CI \\
\midrule
GPT-5.4 & 160 & 0.0608 & [0.0208, 0.1014] & 154 & 0.0610 & [0.0198, 0.1028] \\
GPT-5.4 Mini & 160 & 0.0354 & [-0.0028, 0.0747] & 152 & 0.0241 & [-0.0157, 0.0636] \\
DeepSeek v4 & 160 & 0.1345 & [0.0852, 0.1839] & 147 & 0.1413 & [0.0899, 0.1927] \\
DeepSeek v4 (Thinking) & 159 & 0.1705 & [0.1300, 0.2130] & 141 & 0.1734 & [0.1300, 0.2183] \\
Llama 3.1 70B & 160 & 0.0975 & [0.0606, 0.1355] & 140 & 0.0897 & [0.0512, 0.1302] \\
Llama 3.1 8B & 159 & 0.0851 & [0.0503, 0.1209] & 132 & 0.0715 & [0.0349, 0.1086] \\
Qwen3 14B & 160 & 0.1137 & [0.0729, 0.1556] & 150 & 0.1174 & [0.0748, 0.1615] \\
Qwen3 14B (Thinking) & 160 & 0.1283 & [0.0849, 0.1714] & 145 & 0.1192 & [0.0739, 0.1644] \\
Qwen3 32B & 160 & 0.1281 & [0.0939, 0.1632] & 145 & 0.1345 & [0.0985, 0.1705] \\
Qwen3 32B (Thinking) & 160 & 0.0889 & [0.0563, 0.1217] & 148 & 0.0905 & [0.0567, 0.1246] \\
Qwen3 8B & 160 & 0.0922 & [0.0524, 0.1332] & 142 & 0.0904 & [0.0485, 0.1330] \\
Qwen3 8B (Thinking) & 160 & 0.1326 & [0.0927, 0.1734] & 143 & 0.1321 & [0.0905, 0.1760] \\
\bottomrule
\end{tabular}
\caption{Shared-pair ordering and the subset retaining all six facts in both responses. Intervals are individual 95\% percentile bootstrap intervals from 20,000 scenario resamples. Main effects use the 60-test adjustment; the all-six check uses a separate 12-test adjustment. Eleven models are significant in each analysis; GPT-5.4 Mini is the exception.}
\label{tab:ordering_robustness}
\end{table*}


\section{Neutral Grounding-Instruction Ablation}
\label{app:grounding_ablation}

The original neutral system message contains the sentence ``Use only the
provided facts and do not add outside information.'' Its purpose was to keep
the neutral summary focused on the supplied evidence and discourage outside
additions. Because this instruction is absent from the goal-conditioned
prompt, we tested its contribution by deleting only this sentence from the
neutral prompt for Qwen3-8B (non-thinking) and Llama-3.1-8B. We generated 160
new neutral responses per model and reused the original goal responses,
sentence matches, and framing judgments unchanged. Scenario, user message,
evidence order, institutional role and audience wording were otherwise
unchanged; we did not add a role to the neutral condition.

This yields 320 newly generated responses and 640 outputs in the two-model
paired analysis. Every new response has sentence matches and complete
framing judgments. The experiment used the cached checkpoint generation
defaults (sampling with temperature 0.6 and a 768-token output cap); none of
the new responses reached that cap. The historical outputs lack complete
runtime and stop metadata, so we do not claim exact reconstruction of every
historical runtime setting or establish their absence of truncation.

Table~\ref{tab:grounding_ablation} reports all five dimensions, not only those
with significant results. Both ordering columns use the revised shared-pair
definition. \textit{Selection}, \textit{emphasis}, \textit{numerical specificity
loss}, and \textit{framing} change little in absolute terms. The estimated
\textit{ordering} effect decreases from 0.0922 to 0.0707 for Qwen and from
0.0851 to 0.0545 for Llama, but remains positive and significant. After BH
adjustment across ten tests (five dimensions $\times$ two models),
\textit{selection}, \textit{emphasis}, \textit{ordering}, and \textit{framing}
are significant for both models; \textit{numerical specificity loss} is not
significant for either. These tests assess goal-minus-new-neutral effects,
not changes relative to the original experiment; they are not equivalence
tests. Table counts vary by metric: ordering requires a shared pair,
emphasis requires a nonzero eligible token denominator, and numerical
specificity requires recognized source numbers.

Thus, the main pattern persists in both tested models without the asymmetric
grounding instruction. The experiment does not separate the institutional
role from the objective. The shared user message already supplies recipient
context, which the goal-conditioned system message restates. Roles are
deliberately absent from the neutral baseline because they may imply an
institutional interest even without an explicit objective. JANUS evaluates
this combined institutional conditioning, not the incremental effect of an
objective sentence under an otherwise identical institutional role.

\begin{table*}[t]
\centering
\small
\setlength{\tabcolsep}{3pt}
\begin{tabular}{llrrlrr}
\toprule
Model & Dimension & Original & Sentence removed & 95\% CI (removed) & $n$ & BH $q$ \\
\midrule
Qwen3-8B & Selection & 0.0146 & 0.0229 & [0.0062, 0.0417] & 160 & 0.0242 \\
Qwen3-8B & Emphasis & 0.0788 & 0.0736 & [0.0292, 0.1195] & 158 & 0.0033 \\
Qwen3-8B & Ordering & 0.0922 & 0.0707 & [0.0267, 0.1165] & 160 & 0.0034 \\
Qwen3-8B & Numerical specificity loss & -0.0046 & -0.0055 & [-0.0335, 0.0204] & 132 & 0.7029 \\
Qwen3-8B & Framing & 0.1848 & 0.1799 & [0.1356, 0.2254] & 160 & $<0.0001$ \\
Llama-3.1-8B & Selection & 0.0312 & 0.0312 & [0.0104, 0.0542] & 160 & 0.0147 \\
Llama-3.1-8B & Emphasis & 0.1457 & 0.1499 & [0.1095, 0.1902] & 160 & $<0.0001$ \\
Llama-3.1-8B & Ordering & 0.0851 & 0.0545 & [0.0203, 0.0882] & 159 & 0.0034 \\
Llama-3.1-8B & Numerical specificity loss & 0.0269 & 0.0243 & [-0.0031, 0.0526] & 132 & 0.0990 \\
Llama-3.1-8B & Framing & 0.2342 & 0.2373 & [0.1883, 0.2865] & 160 & $<0.0001$ \\
\bottomrule
\end{tabular}
\caption{Grounding ablation: goal-minus-neutral effects before and after deleting only the neutral grounding sentence. Both ordering columns use the shared-pair metric. The original goal responses are reused. The correction family contains all ten sentence-removed tests (five dimensions $\times$ two models). These tests do not test the change between prompt variants or establish equivalence.}
\label{tab:grounding_ablation}
\end{table*}


\section{Numerical Specificity and Descriptive Aggregation}
\label{app:numerical_specificity}

\textit{Numerical specificity} measures retention of normalized source
numbers. Its complementary loss score is a nondirectional fidelity
diagnostic, not a measure of all qualifications or a judgment that a change
necessarily favors the goal. Numbers were chosen because their retention
can be measured consistently and reproducibly. The score is defined for
132 scenarios with recognized source numbers; non-numerical qualifications,
such as uncertainty or eligibility restrictions, are outside its scope.
Matching numerical values does not verify their semantic attribution.

The direction of a wording change cannot be inferred from the polarity of
the fact whose number disappears. Omitting the magnitude of a cost can
downplay a risk; replacing ``returns increased by 5\%'' with ``returns
increased dramatically'' can exaggerate a benefit. In either case the reader
loses information for judging magnitude. We therefore measure numerical
loss symmetrically for favorable and adverse facts, while \textit{framing}
assesses goal-relative characterization. Numerical loss alone is not proof
of misleadingness: appropriate rounding or compression may preserve the
decision-relevant meaning.

\subsection{Aggregation and figure interpretation}
\label{app:metric_aggregation}

The main table's Average preserves the existing descriptive convention:
\[
A(M)=\frac15\sum_{a\in\mathcal A}\max(0,\overline D_a(M)),
\]
where $\mathcal A=\{\mathrm{sel},\mathrm{emp},\mathrm{ord},\mathrm{spec},
\mathrm{frame}\}$ and $\overline D_a$ is the mean over eligible scenarios for
that dimension. It is not the signed mean of the five columns. Because the
dimensions have different scales and numerical loss is nondirectional,
this summary is not a calibrated measure of misleadingness; dimension-level
findings remain primary.

To preserve the original figures' transformations, category composites
first divide the five category-level means by fixed scales $(2,2,1,1,2)$,
clip to $[-1,1]$, and average their positive parts, multiplied by 100.
The category radar divides these values by the largest observed
model--category composite, so its radii are relative to that maximum.
This normalized plotting composite differs from the main table's raw-scale
Average. Replacing ordering can therefore change both a category's
composite and the radar's common radial normalization.

The condition-profile figure averages positive parts of per-response
selection, emphasis, and framing scores, and numerical specificity loss.
Its ordering points are instead the mean neutral and goal arm scores over
the same shared pairs for each scenario. Their difference equals the
reported mean shared-pair ordering effect. They depend on both responses
and must not be interpreted as independent absolute output scores. Each
dimension has its own displayed axis range. The short image label
``Specificity'' refers to numerical specificity loss.

Human-validation pair scores preserve the original sample-selection
aggregation: normalize each scenario's five effects by $(2,2,1,1,2)$, clip
to $[-1,1]$, and average positive parts on a 0--100 scale. The historical
implementation treats an unavailable component as zero for this descriptive
composite, not for dimension-level inference. We retain the fixed annotated
sample and this aggregation, replacing only ordering when recomputing its
association with human judgments.
